% TOP_PROBLEM: {"id": 5500003, "title": "Voronoi Diagram of Lines in 3D", "category_id": 6, "category": {"id": 6, "name": "geometry", "display_name": "Geometry", "description": "Euclidean and non-Euclidean geometry, geometric structures.", "slug": "geometry", "order_index": 6, "created_at": "2026-07-31T15:26:25.671Z"}, "status": "open", "problem_number": "AMR-054-0003", "set_id": 13}
% FIRST_POSTED: 2026-10-01
% ENTRY_KIND: statement_only
% UnsolvedMath ID 5500003; source catalogue code AMR-054-0003
% !TeX program = lualatex
% Definition and sources only; this file is not an LLM solution attempt.
\documentclass[11pt,a4paper]{article}
\usepackage[margin=25mm]{geometry}
\usepackage{fontspec}
\setmainfont{lmroman10-regular.otf}[BoldFont=lmroman10-bold.otf,
 ItalicFont=lmroman10-italic.otf,BoldItalicFont=lmroman10-bolditalic.otf]
\usepackage{amsmath,amssymb,mathtools}
\usepackage{unicode-math}
\setmathfont{latinmodern-math.otf}
\AtBeginDocument{\let\setminus\smallsetminus}
\usepackage{xurl,hyperref}
\hypersetup{colorlinks=true,urlcolor=blue}
\setcounter{secnumdepth}{0}
\setlength{\parindent}{0pt}
\setlength{\parskip}{5pt}
\setlength{\emergencystretch}{3em}
\begin{document}
\section*{Voronoi Diagram of Lines in 3D}\label{5500003}
{\small UnsolvedMath ID 5500003 · AMR-054-0003 · Geometry · AMR Open Problem Lists}
\subsection{Definitions and mathematical statement}
Let $\mathcal S=\{s_1,\ldots,s_n\}$ be straight lines in $\mathbb R^3$,
or, in the segment version, closed nondegenerate line segments. For
$x\in\mathbb R^3$ put $d(x,s)=\inf_{y\in s}\|x-y\|_2$. The Voronoi cell
of site $s_i$ is
\[
 V_i=\{x:d(x,s_i)\le d(x,s_j)\text{ for every }j\}.
\]
The diagram is the subdivision by nearest-site sets. Its combinatorial
complexity counts connected faces of dimensions zero through three,
including disconnected pieces with the same nearest-site label. One
may use inputs in general position so that coincident bisectors and
nongeneric higher-order ties do not make the count ambiguous.

Determine the largest possible complexity as a function of $n$ for each
of the line and segment versions. The principal conjectural target is
near-quadratic complexity: for every fixed $\varepsilon>0$, an upper bound
$O_\varepsilon(n^{2+\varepsilon})$, together with sharp lower constructions.
The distance is Euclidean distance to the whole site, not distance to a
chosen endpoint or to an intersection with a fixed plane. Bounds only
for a single distance level set do not automatically bound the full
three-dimensional diagram.
\subsection{Short English statement}
How complicated can the nearest-line or nearest-segment subdivision of
three-dimensional space be? Is its worst-case size nearly quadratic?
\subsection{Sources}
\begin{itemize}
\item[S1] ULAM.AI and the UnsolvedMath Contributors, supplied problems.json, record 5500003 (AMR-054-0003), AMR Open Problem Lists. Catalogue identity and source transcription; CC BY 4.0.
\url{https://huggingface.co/datasets/ulamai/UnsolvedMath}.
\item[S2] Open Problems Project - Computational Geometry, Problem 3. Original source indexed by AMR-054-0003.
\url{https://topp.openproblem.net/p3}.
\end{itemize}
\end{document}
